% ============================================================
\section{Conclusion}
\label{sec:conclusion}
% ============================================================

This paper designs, derives, and evaluates a fully automated, rule-based tactical
asset allocation strategy. It rotates monthly across 15 diversified ETFs, with no
leverage, no derivatives, and no discretionary intervention. Four components make it
up: a multi-period (13612U) momentum score, a top-7 breadth pre-filter, an
absolute-momentum threshold that drives a breadth-based defensive switch, and a
minimum-correlation triplet selection rule. Every parameter is derived on the
2004--2015 training window, confirmed (for the triplet size) on a 2016--2020
validation window, and evaluated once on a locked 2021--2025 test block.

The headline numbers are competitive. Over the full 2004--2025 period the strategy
attains a CAGR of 11.00\%, a Sharpe ratio of 0.892, and a maximum drawdown of
$-11.86\%$. On the locked test block it attains 11.42\%, 0.873, and $-7.06\%$: the
highest Sharpe of the benchmark panel, at roughly one-third the drawdown of $1/N$,
60/40, and the S\&P~500. Against naive $1/N$ on the test the edge is marginal on the
Sharpe difference ($p = 0.079$) but robust on a Newey--West spanning alpha
($+6.88\%$/yr, $p = 0.008$), which also holds against the CAPM and a two-factor
equity-plus-duration model. The strategy also beats a Faber (2007) trend rule
($p = 0.032$), so the edge is not merely trend-following. Against 60/40 and the
S\&P~500 it matches Sharpe while controlling drawdown far more tightly.

Three contributions follow. The first gives minimum-correlation triplet selection a
formal optimality footing (Proposition~\ref{prop:mincorr}), then shows empirically
that the rule's value is risk control: it cuts variance and tail risk at
statistically equal expected return, and we present it as such rather than as a
return enhancement. The second is a transparent, two-pronged defence against data
snooping. One prong is a matched-permutation test, in which the momentum rule beats
chance on the Sharpe ratio and the minimum-correlation rule beats chance through
variance reduction; the other is a nested incremental spanning-alpha ladder that
attributes the edge to stacking complementary rules rather than to any single one.
The third is robustness where it matters. The outperformance survives a joint
control for equity, duration, and a generic trend factor (residual alpha
$\approx +5\%$/yr, so it is not merely trend-following), realistic transaction costs
(break-even one-way cost near $62$ bps against a measured turnover of $10\times$ per
year), both halves of the sample, a Deflated Sharpe ratio of $0.975$, and, most
directly, 100\% real (non-proxy) ETF data.

We are explicit about what the strategy does \emph{not} establish. It does not beat
a 60/40 portfolio on a risk-adjusted basis; its advantage over $1/N$ is marginal on
the Sharpe difference and rests more securely on the spanning alpha; and its
moderate-to-high turnover makes net performance more cost-sensitive than
low-turnover peers, though a no-trade-band variant reduces turnover while preserving
the alpha. These are stated openly because the strategy's credibility rests on
honest accounting, not on the headline.

Several extensions follow naturally. A no-trade-band implementation widens the cost
margin, whereas calendar rebalancing does not, because it worsens drawdown. An
expanding-window proxy recalibration would formally close the residual look-ahead
concern, already shown immaterial on real data. A prospective live implementation
would supply the only fully out-of-sample evidence free of the researcher degrees of
freedom that attach to any backtest. A formal replication on UCITS-listed
equivalents would test whether the edge survives the instrument substitution
required for European retail investors; a preliminary prototype on 14 LSE-listed
equivalents (REM excluded; SCZ replaced by WOSC.L, IEF by IDTM.L) already replicates
the post-training spanning alpha ($\alpha = +6.81\%$/yr, $p = 0.0006$ over
2016--2025). The core message is narrow and defensible. A simple, auditable,
four-rule allocation strategy earns a spanning alpha over a naive diversified
benchmark, sourced from momentum and risk-controlled by minimum-correlation
construction, while keeping drawdowns well inside those of passive multi-asset
portfolios.
